\documentclass[11pt, a4paper]{article}

\usepackage[a4paper, margin=2.5cm]{geometry}
\usepackage[utf8]{inputenc}
\usepackage[IL2]{fontenc}
\usepackage[slovak]{babel}
\babelprovide[import, onchar=ids fonts]{slovak}

\usepackage{enumitem}
\setlist[itemize]{label=-}
\usepackage{amsmath}
\usepackage{booktabs}
\usepackage{url}
\usepackage{indentfirst}
\usepackage{hyperref}

\hypersetup{
    colorlinks=true,
    linkcolor=blue,
    citecolor=blue,
    urlcolor=cyan
}

\title{Akustická prediktívna údržba priemyselných zariadení pomocou okrajového strojového učenia a IoT\thanks{Semestrálny projekt v predmete Metódy inžinierskej práce, ak. rok 2026/2027, vedenie: Ing. Ivan Kapustík}}

\author{Gabriela Tovarňáková, Pavel Ustimenko, Ján Vargovčík\\[2pt]
	{\small Slovenská technická univerzita v Bratislave}\\
	{\small Fakulta informatiky a informačných technológií}\\
}

\date{\small 4. október 2026}

\begin{document}

\maketitle



\section{Spresnenie zamerania projektu}

Na základe požiadaviek výzvy pre aplikovaný výskum uvádzame spresnenie troch kľúčových položiek formulára žiadosti:

\subsection*{Názov projektu / Project Title}
\begin{itemize}
    \item \textbf{Slovenský názov:} Akustická prediktívna údržba priemyselných zariadení pomocou okrajového strojového učenia a IoT
    \item \textbf{Anglický názov:} Acoustic Predictive Maintenance of Industrial Equipment Using Edge Machine Learning and IoT
\end{itemize}

\subsection*{Akronym projektu}
\begin{itemize}
    \item \textbf{Akronym:} SOUND-EDGE
\end{itemize}

\subsection*{Anotácia projektu}
Projekt SOUND-EDGE sa zameriava na výskum a vývoj inovatívneho systému prediktívnej údržby pre rotujúce stroje (motory, ložiská) v prostredí Priemyslu 4.0. Tradičné prístupy vyžadujú odosielanie veľkého množstva surových senzorových dát (napríklad vibrácií a akustických signálov) do cloudových úložísk. Tento prístup spôsobuje vysokú latenciu, enormnú spotrebu sieťového pásma a potenciálne bezpečnostné riziká. Komerčné riešenia sú navyše pre mnohé malé a stredné podniky finančne nákladné.

Tento projekt rieši spomenuté problémy využitím technológie TinyML (okrajové strojové učenie -- Edge AI). Hlavným cieľom je navrhnúť architektúru a optimalizované softvérové modely hlbokých neurónových sietí, ktoré sú prispôsobené pre analýzu akustických a vibračných signálov priamo v prostredí nízkonákladových, energeticky obmedzených IoT zariadení. Spracovanie dát priamo na okraji siete výrazne znižuje spotrebu energie spojenú s bezdrôtovým prenosom dát, znižuje latenciu a zvyšuje celkovú bezpečnosť systému~\cite{10781407}. Cloudové služby budú v navrhovanej architektúre využité výhradne na agregáciu sumárnych štatistík anomálií, čím sa minimalizuje objem prenášaných dát.

Hlavným výstupom projektu bude návrh softvérovej architektúry a simulačný model TinyML algoritmu, ktorý bude laboratórne otestovaný a vyhodnotený na dostupných benchmarkových priemyselných datasetoch. Pre overenie vlastností navrhovaného algoritmického riešenia budú v rámci simulačných experimentov použité štandardné súbory dát, predovšetkým dataset NASA Bearing~\cite{NASA:Bearing}, ktorý obsahuje kontinuálne záznamy vibrácií až do štádia zlyhania ložiska. Výskum zabezpečí posun z teoretického konceptu do fázy softvérovo overeného riešenia (posun z TRL 3 na TRL 4). Navrhnutá metodika a simulačné experimenty preukážu schopnosť systému včasne detegovať mikroskopické poruchy, čo podnikom poskytne overený základ pre znižovanie nákladov na údržbu a predlžovanie životnosti výrobných zariadení.\bibliographystyle{unsrt}
\bibliography{literatura}

\end{document}
